\thispagestyle{empty}
\vspace*{8mm}
{\sffamily\bfseries\color{Teal} MAE5920\enspace $\cdot$\enspace INTRODUÇÃO A DEEP LEARNING\par}
\vspace{5mm}
{\fontsize{34}{37}\selectfont\sffamily\bfseries O que a CNN lê\\ em uma resenha?\par}
\vspace{3mm}
{\large Classificação de sentimento e tokenização no IMDb\par}
\vspace{5mm}
{\sffamily Nestor, Letícia e Gabriel\par}
\vspace{6mm}
\noindent\textcolor{Rule}{\rule{\linewidth}{.7pt}}

\subsection*{Resumo}
Classificamos resenhas do IMDb em positivas e negativas com regressão logística
sobre TF-IDF, uma CNN com embeddings aprendidos e encoders pré-treinados.
O experimento original avalia os modelos no teste oficial. Ensaios internos
com três sementes separam escolhas de otimizador, arquitetura e tokenização.
Na comparação pareada, BPE reduziu tokens desconhecidos de 3,27\% para zero e
a acurácia média passou de 82,00\% para 81,41\%. Em um par sintético com as mesmas
palavras, as duas CNNs mudaram de classe quando trocamos a ordem dos adjetivos.
Os seis checkpoints inverteram a interpretação prevista para as duas frases.
Um controle com as mesmas janelas convolucionais produziu logits idênticos
em cada checkpoint. Esses exemplos ajudam a localizar o que a arquitetura
consegue distinguir, mesmo quando o desempenho agregado parece semelhante.

\section{Ordem e sentimento}
\label{sec:ordem}
Considere duas resenhas que só trocam os lugares de \textit{terrible} e
\textit{wonderful}. A oração após \textit{but} expressa a avaliação final.

\vspace{2mm}
\caso{A}{\TextoA}
\caso{B}{\TextoB}
\vspace{1mm}
As duas sequências têm a mesma contagem de cada palavra. Um modelo que recebe
apenas unigramas vê o mesmo vetor para A e B. Os bigramas mudam, portanto o
TF-IDF de referência deste trabalho, que inclui bigramas, pode separá-las.
A CNN também pode responder de modo diferente, pois aplica filtros a janelas
ordenadas de embeddings. A questão empírica é se ela atribui o sentimento à
oração que encerra o contraste.

Investigamos essa questão em três escalas. Primeiro, medimos desempenho em
resenhas reservadas. Depois, alteramos componentes da CNN sob controles de
treino e avaliamos palavras contra subpalavras. Por fim, consultamos pesos
congelados com pares sintéticos para verificar sensibilidade à ordem e
construir um limite exato da arquitetura.

\vfill
\nota{Código, notebook executado e resultados por execução.
\href{\projetourl}{Repositório do projeto}.}
